\section{Methodik}

Verglichen werden fuenf Verfahren: Trivialklassifikator (Baseline),
logistische Regression, Random Forest, Gradient Boosting und ein neuronales
Netz -- jeweils ohne und mit Ueberabtastung der Minderheitsklasse (SMOTE).

\paragraph{Verschachtelte Kreuzvalidierung.}
Die Bewertung laeuft ueber 10 aeussere Folds, die Hyperparameterwahl in einem
inneren 3-Fold-Schritt. Damit sieht die Parameterwahl die Bewertungsdaten nie:
Die Kennzahlen sind Messungen, keine Obergrenzen. SMOTE sitzt in der Pipeline
und wirkt nur auf die Trainingsfolds, damit keine kuenstlichen Punkte in die
Bewertung gelangen.

\paragraph{Gepaarter Signifikanztest.}
Die Fold-Ergebnisse je Modell beruhen auf denselben aeusseren Aufteilungen und
sind damit gepaart. Ein Wilcoxon-Vorzeichen-Rang-Test prueft, ob ein
Unterschied groesser ist als die Streuung zwischen den Folds -- strenger als
der optische Vergleich von Fehlerbalken.

\paragraph{Kostenmodell.}
Zwei Fehler sind moeglich: ein verpasster Ausfall (falsch negativ, teuer) und
ein Fehlalarm (falsch positiv, kostet eine Pruefung). Mit einem
Kostenverhaeltnis von $k{:}1$ laesst sich die Schwelle aus dem erwarteten
Nutzen ableiten, statt sie bei 0,5 festzulegen. \textbf{Das vom Auftraggeber
genannte Verhaeltnis betraegt 50:1} -- ein verpasster Ausfall kostet also das
50-fache einer Pruefung. Dieses Leitverhaeltnis steht im Mittelpunkt; die
Verhaeltnisse 1:1, 5:1 und 10:1 bleiben als Vergleichsszenarien daneben
stehen.
